%! Unit 2: Quadratic Functions — Simplifying Radicals Quiz (KEY)
% GENERATED BODY: the blank and key bodies are byte-identical except the header title;
% answers live in \slot / work, which the preamble shows or hides.
% A cheat sheet leads the front: simplifying with a prime factor tree, and pulling out i
% from a negative radicand (as taught in 2.3). Then two real radicands (one with a
% coefficient) and three negative radicands (perfect square, non-perfect, a + sqrt(-b)).
% Front/back: cheat sheet + 1--2 on the front, 3--5 on the back. Numbers are distinct from
% the cheat-sheet examples, the 2.3 deck, and the unit tests.
\documentclass[12pt]{article}
\usepackage{algebra2-article}
\usepackage{algebra2-key}   % pulls in -boxes; do NOT also load -boxes
\usepackage{grading}
\gradingassignment{alg2-u02-radicals-quiz}
\gradingcourse{algebra_2}
\usepackage{algebra2-graded}
% Question line: forest label, then the prompt on the same line, over a thin rule
\newcommand{\qhead}[2]{%
  \boxguard[8]\vspace{10pt}\noindent
  {\color{forest}\bfseries #1}\hfill #2\par
  \vspace{-2pt}\noindent{\color{forest}\rule{\linewidth}{0.6pt}}\par}
% THE WORK RULE for a test: \slot{width}{answer} is the ONLY line that differs between
% blank and key, and it differs in the preamble, not the body.
\newcommand{\slot}[2]{\underline{\makebox[#1]{\ans{#2}}}}
% Cheat-sheet factor tree node styles
\tikzset{ftree/.style={font=\small, inner sep=1.5pt},
         fprime/.style={ftree, circle, draw=forest, thick, minimum size=13pt}}

\begin{document}

\pageheader{Unit 2: Quadratic Functions}{Simplifying Radicals Quiz --- Answer Key}

\vspace{0.04in}
\noindent\begin{minipage}[t]{0.55\linewidth}
\vspace{0pt}
\begin{fixedskillbox}[Cheat Sheet: Prime Factor Tree]{forestbg}
\small
\begin{minipage}[t]{0.40\linewidth}
\vspace{2pt}
\begin{tikzpicture}[level distance=8.5mm,
    level 1/.style={sibling distance=13mm},
    level 2/.style={sibling distance=13mm},
    level 3/.style={sibling distance=6.5mm}]
  \node[ftree] {$72$}
    child {node[fprime] {$2$}}
    child {node[ftree] {$36$}
      child {node[ftree] {$6$}
        child {node[fprime] {$2$}}
        child {node[fprime] {$3$}}}
      child {node[ftree] {$6$}
        child {node[fprime] {$2$}}
        child {node[fprime] {$3$}}}};
\end{tikzpicture}
\end{minipage}\hfill
\begin{minipage}[t]{0.57\linewidth}
\vspace{2pt}
\textbf{1.} Break the radicand down until every branch ends in a \textbf{prime} (the circled numbers).\\[2pt]
\textbf{2.} Group matching primes into \textbf{pairs}.\\[2pt]
\textbf{3.} Each pair leaves the root as \textbf{one} number. Leftovers stay inside.\\[2pt]
\textbf{4.} Multiply what came out by any number already in front.
\end{minipage}

\vspace{4pt}
$\sqrt{72}=\sqrt{(2\cdot2)(3\cdot3)\cdot2}=2\cdot3\sqrt{2}=6\sqrt{2}$
\end{fixedskillbox}
\end{minipage}\hfill
\begin{minipage}[t]{0.43\linewidth}
\vspace{0pt}
\begin{fixedskillbox}[Cheat Sheet: Negative Radicand]{goldbg}
\small
\[ \sqrt{-1}=i \qquad\text{so}\qquad \sqrt{-k}=i\sqrt{k} \]
\textbf{1.} Pull the $i$ out \textbf{first}.\\[2pt]
\textbf{2.} Simplify $\sqrt{k}$ with a factor tree.\\[2pt]
\textbf{3.} Write the $i$ \textbf{before} the radical: $2i\sqrt{3}$, not $2\sqrt{3}\,i$.

\vspace{4pt}
$\sqrt{-12}=i\sqrt{12}=i\sqrt{2\cdot2\cdot3}=2i\sqrt{3}$

\vspace{3pt}
\textbf{Watch out:} $\sqrt{-12}\neq-\sqrt{12}$. A negative \emph{inside} the root makes $i$, not a minus sign.
\end{fixedskillbox}
\end{minipage}

\vspace{0.08in}
\noindent\textbf{Instructions:} No calculator. Show your factor tree or your work. Simplify
\textbf{completely}. Write each final answer in its box. Only what is inside a box is graded.
\hfill\textbf{20 pts each}

% ======================================================================
\qhead{1.\ Simplify the Radical}{Simplify: \quad $\sqrt{98}$}
\begin{work}
  98 &= 2\cdot49 = 2\cdot\underbrace{7\cdot7}_{\text{pair}} \\
  \sqrt{98} &= 7\sqrt{2}
\end{work}
\vfill
\noindent\hfill Answer: \inblank{\slot{4.2cm}{$7\sqrt{2}$}}{\gradedblank[id=q1, type=short, points=20, answer={7sqrt(2)}, check=expression, display={$7\sqrt{2}$}]{4.2cm-2pt}{18pt}}

% ======================================================================
\qhead{2.\ Simplify the Radical}{Simplify: \quad $3\sqrt{80}$}
\begin{work}
  80 &= 8\cdot10 = \underbrace{2\cdot2}_{\text{pair}}\cdot\underbrace{2\cdot2}_{\text{pair}}\cdot5 \\
  \sqrt{80} &= 2\cdot2\sqrt{5} = 4\sqrt{5} \\
  3\sqrt{80} &= 3\cdot4\sqrt{5} = 12\sqrt{5}
\end{work}
\vfill
\noindent\hfill Answer: \inblank{\slot{4.2cm}{$12\sqrt{5}$}}{\gradedblank[id=q2, type=short, points=20, answer={12sqrt(5)}, check=expression, display={$12\sqrt{5}$}]{4.2cm-2pt}{18pt}}

% ======================================================================
\newpage
\qhead{3.\ Negative Radicand}{Simplify: \quad $\sqrt{-81}$}
\begin{work}
  \sqrt{-81} &= i\sqrt{81} \\
  &= 9i
\end{work}
\vfill
\noindent\hfill Answer: \inblank{\slot{4.2cm}{$9i$}}{\gradedblank[id=q3, type=short, points=20, answer={9i}, check=expression]{4.2cm-2pt}{18pt}}

% ======================================================================
\qhead{4.\ Negative Radicand}{Simplify: \quad $\sqrt{-45}$}
\begin{work}
  \sqrt{-45} &= i\sqrt{45} \\
  45 &= 5\cdot9 = 5\cdot\underbrace{3\cdot3}_{\text{pair}} \\
  i\sqrt{45} &= 3i\sqrt{5}
\end{work}
\vfill
\noindent\hfill Answer: \inblank{\slot{4.2cm}{$3i\sqrt{5}$}}{\gradedblank[id=q4, type=short, points=20, answer={3sqrt(5)i}, check=expression, display={$3i\sqrt{5}$}]{4.2cm-2pt}{18pt}}

% ======================================================================
\qhead{5.\ Negative Radicand in $a+bi$ Form}{Simplify: \quad $6+\sqrt{-28}$}
\begin{work}
  \sqrt{-28} &= i\sqrt{28} \\
  28 &= 7\cdot4 = 7\cdot\underbrace{2\cdot2}_{\text{pair}} \\
  i\sqrt{28} &= 2i\sqrt{7} \\
  6+\sqrt{-28} &= 6+2i\sqrt{7}
\end{work}
\vfill
\noindent\hfill Answer: \inblank{\slot{4.2cm}{$6+2i\sqrt{7}$}}{\gradedblank[id=q5, type=short, points=20, answer={6+2sqrt(7)i}, check=expression, display={$6+2i\sqrt{7}$}]{4.2cm-2pt}{18pt}}

\end{document}
